\documentclass[border=10pt]{standalone}
\usepackage{tikz,ctex}
\usetikzlibrary{positioning}

\begin{document}

\begin{tikzpicture}[
    % 全局设置
    node distance = 1.6cm and 1.4cm,
    every node/.style = {
        draw,
        rounded corners,
        thick,
        align = center,
        font = \sffamily\small,
        minimum width = 2.0cm,
        minimum height = 0.9cm
    },
    % 样式定义
    centerStyle/.style = {
        fill = blue!45,
        text = white,
        font = \sffamily\bfseries\large,
        line width = 2pt,
        minimum size = 2.5cm,
        circle
    },
    branchTitle/.style = {
        font = \sffamily\bfseries,
        line width = 1.5pt
    },
    subNode/.style = {
        fill = white,
        font = \sffamily\small,
        minimum width = 2.0cm,
        minimum height = 0.8cm
    },
    leafNode/.style = {
        fill = white,
        font = \sffamily\footnotesize,
        minimum width = 1.8cm,
        minimum height = 0.7cm
    },
    arrowStyle/.style = {
        ->,
        thick,
        >=stealth,
        shorten >=2pt
    }
]

    % ================= 1. 中心节点 =================
    \node[centerStyle] (center) {残差连接 \\ Residual Connections};

    % ================= 2. 上方分支：深度网络问题 (红) =================
    \node[branchTitle, fill=red!15, above=of center] (problem) {深度网络问题 \\ Deep Network Issues};

    \node[subNode, above left=of problem] (vanish) {梯度消失/爆炸 \\ Vanishing/Exploding};
    \node[subNode, above=of problem] (shatter) {破碎梯度 \\ Shattered Gradients};
    \node[subNode, above right=of problem] (deep) {深度训练困难 \\ Hard to Train Deep};
    \node[subNode, right=of problem, yshift=0cm] (bn) {批归一化 \\ Batch Normalization};

    % ================= 3. 右侧分支：残差块结构 (蓝) =================
    \node[branchTitle, fill=orange!15, right=of center] (block) {残差块结构 \\ Residual Block};

    \node[subNode, above=of block] (identity) {恒等映射 \\ Identity Mapping};
    \node[subNode, above right=of block] (skip) {跳跃连接 \\ Skip Connection};
    \node[subNode, right=of block] (residual) {学习残差 \\ Learn Residual};
    \node[subNode, below right=of block, yshift=0cm] (dim) {维度匹配 \\ Dimension Match};
    \node[subNode, below=of block] (nonlin) {非线性函数 \\ Nonlinear $\mathbf{F}_i$};

    % ================= 4. 下方分支：残差网络特性 (绿) =================
    \node[branchTitle, fill=green!15, below=of center] (resnet) {残差网络特性 \\ ResNet Properties};

    \node[subNode, below left=of resnet] (smooth) {平滑误差面 \\ Smooth Error Surface};
    \node[subNode, below=of resnet] (parallel) {多子网络并行 \\ Parallel Sub-networks};
    \node[subNode, below right=of resnet] (hundreds) {数百层可训练 \\ Hundreds of Layers};
    \node[subNode, left=of resnet, yshift=0cm] (grad) {梯度更稳定 \\ Stable Gradients};

    % ================= 5. 左侧分支：实现细节 (紫) =================
    \node[branchTitle, fill=purple!15, left=of center] (impl) {实现细节 \\ Implementation};

    \node[subNode, above left=of impl] (order) {连接位置 \\ Connection Order};
    \node[subNode, left=of impl] (relu) {ReLU 之后 \\ After ReLU};
    \node[subNode, below left=of impl] (preact) {预激活变体 \\ Pre-activation};
    \node[subNode, above=of impl, yshift=-0.3cm] (batchnorm) {批归一化 \\ Batch Norm};

    % ================= 连线逻辑 =================
    % 中心 -> 四大分支标题
    \draw[arrowStyle, red] (center) -- (problem);
    \draw[arrowStyle, blue] (center) -- (block);
    \draw[arrowStyle, green!60!black] (center) -- (resnet);
    \draw[arrowStyle, purple] (center) -- (impl);

    % 分支标题 -> 子节点 (上方)
    \draw[arrowStyle, red!60] (problem) -- (vanish);
    \draw[arrowStyle, red!60] (problem) -- (shatter);
    \draw[arrowStyle, red!60] (problem) -- (deep);
    \draw[arrowStyle, red!60] (problem) -- (bn);

    % 分支标题 -> 子节点 (右侧)
    \draw[arrowStyle, blue!60] (block) -- (identity);
    \draw[arrowStyle, blue!60] (block) -- (skip);
    \draw[arrowStyle, blue!60] (block) -- (residual);
    \draw[arrowStyle, blue!60] (block) -- (dim);
    \draw[arrowStyle, blue!60] (block) -- (nonlin);

    % 分支标题 -> 子节点 (下方)
    \draw[arrowStyle, green!60!black] (resnet) -- (smooth);
    \draw[arrowStyle, green!60!black] (resnet) -- (parallel);
    \draw[arrowStyle, green!60!black] (resnet) -- (hundreds);
    \draw[arrowStyle, green!60!black] (resnet) -- (grad);

    % 分支标题 -> 子节点 (左侧)
    \draw[arrowStyle, purple!60] (impl) -- (order);
    \draw[arrowStyle, purple!60] (impl) -- (relu);
    \draw[arrowStyle, purple!60] (impl) -- (preact);
    \draw[arrowStyle, purple!60] (impl) -- (batchnorm);

\end{tikzpicture}

\end{document}